% !PRINTED 11
\weberpage{11}

%<<CONT
\gesperrt{If $Q$ is a divisor of $P$ which is identical with all of its conjugate
divisors, then $Q$ is called a normal divisor of
$P$}\footnote{Also called a distinguished or invariant subgroup.}.

The group formed from the single element 1 is a normal divisor of
every group. We obtain, further, a normal divisor in the greatest
common divisor $R$ of all the divisors conjugate to any divisor
$Q$ of $P$.

For it has already been proved above that $R$, as the intersection of
several divisors of $P$, is a group.

If now
\weq{1}{Q,\ Q',\ Q'' \ldots}
are the divisors of $P$ conjugate to $Q$, and $b$ is any element of $P$,
then the system of groups
\weq{2}{b^{-1}Qb,\ b^{-1}Q'b,\ b^{-1}Q''b \ldots}
is not different from the system (1). But if $R$ is the intersection of the
groups (1), then $b^{-1}Rb$ is the intersection of (2), and consequently
$R$ is identical with $b^{-1}Rb$, i.e. $R$ is a normal divisor of $P$.

If $N$ is a normal divisor of any group $P$, and $b$ an arbitrary
element in $P$, then
\weq{3}{b^{-1}Nb = N \quad \text{or} \quad Nb = bN.}

If $P$ is of degree $n$, $N$ of degree $\nu$ and of index $\mu$, so that
$n = \mu\nu$, then we can choose the $\mu$ elements $1, b_1 \ldots b_{\mu-1}$ in
such a way that the decomposition of $P$ into the subsidiary groups
\[
\begin{aligned}
P &= N + Nb_1 + Nb_2 + \cdots + Nb_{\mu-1} \\
  &= N + b_1 N + b_2 N + \cdots + b_{\mu-1} N
\end{aligned}
\]
results. Thus
\weq{4}{N,\ N_1 = Nb_1 = b_1 N,\ N_2 = Nb_2 = b_2 N \ldots,\quad N_{\mu-1} = Nb_{\mu-1} = b_{\mu-1} N}
is the system of the subsidiary groups.

If $a$ is any element in $N$, then $Na$ is identical with $N$; hence
$Nab$ is also identical with $Nb$, and since every element $c$ of $P$
occurs in $N$ or in one of its subsidiary groups, and is therefore
contained in the form $ab_k$, $Nc = cN$ must agree with one of the
systems (4).
